\documentclass[unicode,11pt,a4paper,oneside,numbers=endperiod,openany]{scrartcl}

\usepackage{amsmath}
\usepackage[table,xcdraw]{xcolor}
\usepackage{soul}
\sethlcolor{lightgray}
\usepackage{float}
\usepackage{tikz}
\usetikzlibrary{matrix}
\usepackage{amsfonts}

\usepackage{ifthen}
\usepackage[utf8]{inputenc}
\usepackage{graphics}
\usepackage{graphicx}
\usepackage{hyperref}

\pagestyle{plain}
\voffset -5mm
\oddsidemargin  0mm
\evensidemargin -11mm
\marginparwidth 2cm
\marginparsep 0pt
\topmargin 0mm
\headheight 0pt
\headsep 0pt
\topskip 0pt        
\textheight 255mm
\textwidth 165mm

\newcommand{\duedate} {}
\newcommand{\setduedate}[1]{%
\renewcommand\duedate {Due date:~ #1}}
\newcommand\isassignment {false}
\newcommand{\setassignment}{\renewcommand\isassignment {true}}
\newcommand{\ifassignment}[1]{\ifthenelse{\boolean{\isassignment}}{#1}{}}
\newcommand{\ifnotassignment}[1]{\ifthenelse{\boolean{\isassignment}}{}{#1}}

\newcommand{\assignmentpolicy}{
\begin{table}[h]
\begin{center}
\scalebox{0.8} {%
\begin{tabular}{|p{0.02cm}p{16cm}|}
\hline
&\\
\multicolumn{2}{|c|}{\Large\textbf{HPC Lab for CSE 2021 ---  Submission Instructions}}\\
\multicolumn{2}{|c|}{\large\textbf{(Please, notice that following instructions are mandatory: }}\\
\multicolumn{2}{|c|}{\large\textbf{submissions that don't comply with, won't be considered)}}\\
&\\
\textbullet & Assignments must be submitted to \href{https://moodle-app2.let.ethz.ch/course/view.php?id=14316}{Moodle} (i.e. in electronic format).\\
\textbullet & Provide both executable package and sources (e.g. C/C++ files, Matlab). 
If you are using libraries, please add them in the file. Sources must be organized in directories called:\\
\multicolumn{2}{|c|}{\textit{Project\_number\_lastname\_firstname}}\\
& and  the  file must be called:\\
\multicolumn{2}{|c|}{\textit{project\_number\_lastname\_firstname.zip}}\\
\multicolumn{2}{|c|}{\textit{project\_number\_lastname\_firstname.pdf}}\\
\textbullet &  The TAs will grade your project by reviewing your project write-up, and looking at the implementation 
                 you attempted, and benchmarking your code's performance.\\

\textbullet & You are allowed to discuss all questions with anyone you like; however: (i) your submission must list anyone you discussed problems with and (ii) you must write up your submission independently.\\
\hline
\end{tabular}
}
\end{center}
\end{table}
}
\newcommand{\punkte}[1]{\hspace{1ex}\emph{\mdseries\hfill(#1~\ifcase#1{Points}\or{Points}\else{Points}\fi)}}


\newcommand\serieheader[6]{
\thispagestyle{empty}%
\begin{flushleft}
\includegraphics[width=0.4\textwidth]{eth_logo_kurz_pos_13}
\end{flushleft}
  \noindent%
  {\large\ignorespaces{\textbf{#1}}\hspace{\fill}\ignorespaces{ \textbf{#2}}}\\ \\%
  {\large\ignorespaces #3 \hspace{\fill}\ignorespaces #4}\\
  \noindent%
  \bigskip
  \hrule\par\bigskip\noindent%
  \bigskip {\ignorespaces {\Large{\textbf{#5}}}
  \hspace{\fill}\ignorespaces \large \ifthenelse{\boolean{\isassignment}}{\duedate}{#6}}
  \hrule\par\bigskip\noindent%  \linebreak
 }

\makeatletter
\def\enumerateMod{\ifnum \@enumdepth >3 \@toodeep\else
      \advance\@enumdepth \@ne
      \edef\@enumctr{enum\romannumeral\the\@enumdepth}\list
      {\csname label\@enumctr\endcsname}{\usecounter
        {\@enumctr}%%%? the following differs from "enumerate"
	\topsep0pt%
	\partopsep0pt%
	\itemsep0pt%
	\def\makelabel##1{\hss\llap{##1}}}\fi}
\let\endenumerateMod =\endlist
\makeatother




\usepackage{textcomp}





\begin{document}


\setassignment
\setduedate{May 20, 2021, 12pm (midnight)}

\serieheader{High-Performance Computing Lab for CSE}
    {2021}
    {Student: Pascal Dominic Müller}
    {\\ \\ Discussed with: \\
        $\bullet$ Tarzis Maurer [tamaurer] \\
        $\bullet$ Zuzanna Herud [zherud] \\
    }
    {Solution for Project 6}
    {}
\newline

\assignmentpolicy

% -------------------------------------------------------------------------- %
% -------------------------------------------------------------------------- %
% --- Exercise 1 ----------------------------------------------------------- %
% -------------------------------------------------------------------------- %
% -------------------------------------------------------------------------- %

\section{Scientific Mathematical HPC Software Frameworks - The Poisson Equation [35 points]}

\subsection{Problem Statement: The Poisson Equation}

We solve the poisson equation on the 2 dimensional unit rectangle
$\Omega \subset \mathbb R^2=[0,1] \times [0,1]$ with vanishing dirichlet BC.
\begin{alignat}{3}
  -\Delta &y(x_1, x_2) &&= f(x_1, x_2) &&\text{  on  } \Omega = [0,1] \times [0,1]\\
          &y(x_1, x_2) &&= 0 &&\text{  on  } \partial \Omega
\end{alignat}

whereas $\Delta \equiv \frac{\partial^2}{\partial x_1^2} + \frac{\partial^2}{\partial x_2^2}$

\subsection{Discretization}


\subsubsection{Theory}
Let $g(x,y)$ be an at least twice continuous differentiable function. For the
discretization we use a central finite difference approach [1][Project 3].

The central finite difference in the direction of x (similar for y) is given by:

\begin{equation}
  \delta_h [g](x) = g(x+\frac{1}{2}h) - g(x-h)
\end{equation}

which leads to

\begin{equation}
  \frac{\partial g(x)^2}{\partial^2 x} = \dots = 
    \frac{g(x + h) - 2g(x) + g(x - h)}{h^2}
\end{equation}

\subsubsection{LHS}

Note: here $y$ denotes the seeked function, not the coordinate.

We first write out the laplacian operator.
\begin{equation}
  -\Delta y(x_1, x_2) = \frac{\partial y(x_1, x_2)^2}{\partial x_1^2} 
                      + \frac{\partial y(x_1, x_2)^2}{\partial x_2^2}
\end{equation}

Applying (4) to (5) for $x_1$ and $x_2$ separately, we get:
\begin{equation}
  \frac{\partial y(x_1, x_2)^2}{\partial x_1^2} = 
    \frac{y(x_1 + h, x_2) - 2y(x_1, x_2) + y(x_1 - h, x_2)}{h^2}
\end{equation}
\begin{equation}
  \frac{\partial y(x_1, x_2)^2}{\partial x_1^2} = 
    \frac{y(x_1, x_2 + h) - 2y(x_1, x_2) + y(x_1, x_2 - h)}{h^2}
\end{equation}

We plug (6) and (7) into (5) and use the notations (similar for $x_2$ and $j$
\begin{itemize}
  \item $y(x_1, x_2) = y_{i,j}$
  \item $y(x_1 + h, x_2) = y_{i+1, j}$
  \item $y(x_1 - h, x_2) = y_{i-h, j}$
\end{itemize}

to end up with

\begin{align}
  -\Delta y(x_1,x_2) & \\
                     &= \frac{y(x_1 + h, x_2) - 2y(x_1, x_2) + y(x_1 - h, x_2)}{h^2} \\
                     &\ \ + \frac{y(x_1, x_2 + h) - 2y(x_1, x_2) + y(x_1, x_2 - h)}{h^2} \\
                   &= \frac{1}{h^2} ( y_{i+1,j} - 2y_{i,j} + y_{i-1,j}
                      + y_{i, j+1} - 2y_{i,j} + y_{i,j-1}) \\
                   &= \frac{1}{h^2} ( -4_{i,j} + y_{i+1,j} + y_{i-1,j}
                      + _{i, j+1} + y_{i,j-1} )
\end{align}

We get:

\begin{equation}
  \boxed{
      -\Delta y(x_1,x_2) = \frac{1}{h^2} ( -4y_{i,j} + y_{i+1,j} + y_{i-1,j}
                           + y_{i, j+1} + y_{i,j-1} )
  }
\end{equation}

\subsubsection{RHS}
On the RHS we have a function $f(x_1, x_2): \mathbb R^2 \to \mathbb R$.

Its discretization is given by 

\begin{equation}
  \boxed{
    f_{i,j} := f(x_1^i, x_2^j)
  }
\end{equation}

\subsection{Discretized Problem}


\subsection{Refs}
[1]: https://en.wikipedia.org/wiki/Finite\_difference


\begin{table}[h]
	\caption{Wall-clock time (in seconds) and speed-up (in brackets) using multiple cores on Euler for solving the Poisson problem.}
	\centering
	
	\begin{tabular}{l|r||r|r|r|r}\hline\hline
		Problem & \multicolumn{1}{c||}{$N$} &  \multicolumn{4}{c}{Number of Euler cores} \\
		&       & \multicolumn{1}{c|}{1} & \multicolumn{1}{c|}{8} & \multicolumn{1}{c|}{16} & \multicolumn{1}{c}{32} \\
		\hline\hline
		{ Poisson} & $500^2$  &    \phantom{222222}        &    \phantom{222222}      & \phantom{222222}         &      \phantom{222222} \\
		{ Poisson} & $1000^2$ &            &          &          &       \\
		{ Poisson} & $2000^2$ &            &          &          &       \\
		{ Poisson} & $3000^2$ &            &          &          &       \\\hline \hline
	\end{tabular}
	
	\label{tab:PDEparallel1}
\end{table}


% -------------------------------------------------------------------------- %
% -------------------------------------------------------------------------- %
% --- Exercise 2 ----------------------------------------------------------- %
% -------------------------------------------------------------------------- %
% -------------------------------------------------------------------------- %

\section{Interactive Supercomputing using Jupyter Notebook  [10 points]}


% -------------------------------------------------------------------------- %
% -------------------------------------------------------------------------- %
% --- Exercise 3 ----------------------------------------------------------- %
% -------------------------------------------------------------------------- %
% -------------------------------------------------------------------------- %


\section{Jupyter Notebook - Parallel PDE-Constrained Optimization [40 points]}

\begin{table}[h]
  \caption{Wall-clock time (in seconds) and speed-up (in brackets) using multiple cores on Euler for solving the PDE-constrained optimization problem.}
	\centering
	
	\medskip
	
	%\footnotesize
	\begin{tabular}{l|r||r|r|r|r}\hline\hline
		Problem & \multicolumn{1}{c||}{$N$} &  \multicolumn{4}{c}{Number of Euler cores} \\
		&       & \multicolumn{1}{c|}{1} & \multicolumn{1}{c|}{8} & \multicolumn{1}{c|}{16} & \multicolumn{1}{c}{32} \\
		\hline\hline
		{ Inverse Poisson} & $500^2$  &    \phantom{222222}        &    \phantom{222222}      & \phantom{222222}         &      \phantom{222222} \\
		{ Inverse  Poisson} & $1000^2$ &            &          &          &       \\
		{ Inverse Poisson} & $2000^2$ &            &          &          &       \\
		{ Inverse Poisson} & $3000^2$ &            &          &          &       \\\hline \hline
	\end{tabular}
	
	\label{tab:PDEparallel}
\end{table}

\section{Task:  Quality of the Report [15 Points]}
Each project will have 100 points (out of  which 15 points will be given to the general quality of the written report).


\section*{Additional notes and submission details}
Submit the source code files (together with your used \texttt{Makefile}) in
an archive file (tar, zip, etc.), and summarize your results and the
observations for all exercises by writing an extended Latex report.
Use the Latex template provided on the webpage and upload the Latex summary
as a PDF to \href{https://moodle-app2.let.ethz.ch/course/view.php?id=14316}{Moodle}.

\begin{itemize}
	\item Your submission should be a gzipped tar archive, formatted like project\_number\_lastname\_firstname.zip or project\_number\_lastname\_firstname.tgz. 
	It should contain
	\begin{itemize}
		\item all the source codes of your solutions;
		\item your write-up with your name  project\_number\_lastname\_firstname.pdf.
	\end{itemize}
	\item Submit your .zip/.tgz through Moodle.
\end{itemize}

\end{document}
